%% =============================================================================
%% Pollity.in — मराठी परिचय पत्रिका
%% =============================================================================
\documentclass[11pt, a4paper]{article}

\usepackage{fontspec}
\usepackage{polyglossia}
\setmainlanguage{marathi}
\setotherlanguage{english}

\setmainfont{ITF Devanagari Marathi}[
  BoldFont   = {ITF Devanagari Marathi Demi},
  ItalicFont = {ITF Devanagari Marathi Book},
]
\newfontfamily\englishfont{Helvetica Neue}

\usepackage[margin=2.5cm, top=2cm, bottom=2cm]{geometry}
\usepackage{xcolor}
\usepackage{titlesec}
\usepackage[most]{tcolorbox}
\usepackage{enumitem}
\usepackage{parskip}
\usepackage{microtype}
\usepackage{tabularx}
\usepackage{booktabs}
\usepackage{hyperref}
\usepackage{fancyhdr}
\usepackage{amssymb}

%% ── Colours ──────────────────────────────────────────────────────────────────
\definecolor{pollityblue}{RGB}{31, 97, 141}
\definecolor{pollitygold}{RGB}{212, 172, 13}
\definecolor{lightgray}{RGB}{245, 245, 245}

%% ── Section Formatting ───────────────────────────────────────────────────────
\titleformat{\section}{\large\bfseries\color{pollityblue}}{}{0em}{}[\color{pollityblue}\titlerule]
\titleformat{\subsection}{\normalsize\bfseries\color{pollityblue}}{}{0em}{}

%% ── Header / Footer ──────────────────────────────────────────────────────────
\pagestyle{fancy}
\fancyhf{}
\lhead{\color{pollityblue}\textbf{\textenglish{Pollity.in}}}
\rhead{\color{gray}\small आम्ही काय बनवत आहोत}
\cfoot{\color{gray}\small \textenglish{pollity.in} \textbar{} मार्च २०२६}
\renewcommand{\headrulewidth}{0.4pt}

\hypersetup{colorlinks=true, urlcolor=pollityblue, linkcolor=pollityblue}

%% =============================================================================
\begin{document}

%% ── शीर्षक ───────────────────────────────────────────────────────────────────
\begin{center}
  {\color{pollityblue}\rule{\linewidth}{1.5pt}}\\[0.5em]
  {\Huge\bfseries\color{pollityblue}\textenglish{Pollity.in}}\\[0.3em]
  {\large निवडून आलेल्या लोकप्रतिनिधींसाठी डिजिटल मुख्य सहाय्यक}\\[0.3em]
  {\color{pollityblue}\rule{\linewidth}{1.5pt}}
\end{center}

\vspace{0.8em}

%% ── समस्या ───────────────────────────────────────────────────────────────────
\section{समस्या काय आहे?}

भारतात \textbf{३३ लाखांहून अधिक निवडून आलेले लोकप्रतिनिधी} आहेत: खासदार, आमदार, नगरसेवक, ग्रामपंचायत सरपंच. यांच्यापैकी प्रत्येकजण दररोज शेकडो नागरिकांच्या समस्यांना सामोरे जातो.

या समस्या येतात \textenglish{WhatsApp} संदेशांतून, प्रत्यक्ष भेटींतून, फोन कॉल्सद्वारे आणि पत्रांतून. पण या सगळ्या तक्रारी एकत्र ठेवण्यासाठी कोणतीही व्यवस्था नाही.

याचा परिणाम असा होतो:

\begin{itemize}[leftmargin=1.5em, itemsep=3pt]
  \item तक्रारी नोंदवल्यानंतर विसरल्या जातात किंवा हरवतात
  \item नागरिकांना त्यांच्या तक्रारीबद्दल कोणतेही उत्तर मिळत नाही
  \item कोणत्या समस्या सुटल्या आणि कोणत्या सुटल्या नाहीत याचा कोणताही हिशोब नसतो
  \item कार्यालयातील कर्मचारी संदेश पुढे पाठवण्यात आणि विभागांचा पाठपुरावा करण्यात वेळ घालवतात
\end{itemize}

"कागदावर सोडवली" आणि "प्रत्यक्षात सोडवली" यांमधील दरी प्रचंड मोठी आहे. कोणीही ती मोजत नाही.

%% ── आम्ही काय बनवत आहोत ──────────────────────────────────────────────────────
\section{आम्ही काय बनवत आहोत?}

\begin{tcolorbox}[colback=pollityblue!8, colframe=pollityblue, sharp corners, boxrule=0.6pt]
\textbf{\textenglish{Pollity.in} हे निवडून आलेल्या लोकप्रतिनिधींसाठी एक \textenglish{AI}-आधारित व्यवस्थापन साधन आहे.}
हे एका डिजिटल मुख्य सहाय्यकासारखे काम करते जे मतदारसंघ कार्यालयाचे दैनंदिन काम सांभाळते, जेणेकरून प्रतिनिधी खऱ्या कामावर लक्ष केंद्रित करू शकतात.
\end{tcolorbox}

\vspace{0.8em}

पहिले मॉड्यूल, \textbf{जनश्रवण}, तक्रार व्यवस्थापन हाताळते:

\begin{enumerate}[leftmargin=1.5em, itemsep=4pt]
  \item नागरिक त्यांची समस्या \textenglish{WhatsApp}वर संदेश पाठवून सांगतात
  \item आमची \textenglish{AI} तंत्रज्ञान तत्काळ तक्रारीचा प्रकार (पायाभूत सुविधा, आरोग्य, भ्रष्टाचार इ.) आणि तातडी ठरवते
  \item नागरिकाला काही सेकंदांत एक संदर्भ क्रमांक मिळतो, जेणेकरून त्यांना समजते की तक्रार नोंदली गेली
  \item कार्यालयातील कर्मचाऱ्यांना एकाच \textenglish{Dashboard}वर सर्व तक्रारी दिसतात, प्रत्येकाला जबाबदार व्यक्ती नेमता येते आणि निराकरणापर्यंत पाठपुरावा करता येतो
  \item प्रतिनिधींना \textenglish{data} मिळतो: किती तक्रारी आल्या, किती सुटल्या, आणि मतदारसंघातील मुख्य समस्या काय आहेत
\end{enumerate}

\vspace{0.5em}

\textbf{\textenglish{WhatsApp}वरील संदेश हरवणार नाहीत. नागरिकांना उत्तर न मिळणे संपेल. अनिश्चितता नाही.}

%% ── मोठी दृष्टी ───────────────────────────────────────────────────────────────
\section{मोठी दृष्टी}

जनश्रवण हे पहिले पाऊल आहे. संपूर्ण \textenglish{Pollity} व्यासपीठात हे समाविष्ट असेल:

\begin{tabularx}{\linewidth}{lX}
\toprule
\textbf{मॉड्यूल} & \textbf{कार्य} \\
\midrule
जनश्रवण      & तक्रार व्यवस्थापन आणि नागरिक सेवा \\
विधायक मित्र & विधिमंडळ संशोधन: विधेयकांचे सारांश, धोरण विश्लेषण, चर्चेचे मुद्दे \\
विकास ट्रॅकर & मतदारसंघ विकास निधीचा मागोवा, प्रकल्पांची स्थिती, अर्थसंकल्पाचा वापर \\
नागरिक सेतू  & \textenglish{Social Media} आणि मास कम्युनिकेशनद्वारे नागरिकांशी संपर्क \\
\bottomrule
\end{tabularx}

\vspace{0.5em}
हे सर्व मॉड्यूल्स एकत्र येऊन प्रत्येक निवडून आलेल्या प्रतिनिधीला ती पायाभूत सुविधा देतात जी आजपर्यंत फक्त मोठ्या राजकारण्यांनाच परवडणारी होती, आणि हे अत्यंत कमी खर्चात.

%% ── हे का महत्त्वाचे आहे ─────────────────────────────────────────────────────
\section{हे का महत्त्वाचे आहे?}

भारताचे लोकशाही तंत्र ३३ लाख स्थानिक प्रतिनिधींवर चालते. त्यांच्यापैकी बहुतेकजण पहिल्यांदाच राजकारणात आलेले आहेत, ज्यांच्याकडे कर्मचारी नाहीत, व्यवस्था नाही, \textenglish{data} नाही. \textenglish{Pollity} त्यांना ही व्यवस्था देते.

जेव्हा महाराष्ट्रातील एखाद्या ग्रामीण भागातील सरपंच आपल्या मतदारसंघाला दाखवू शकते की गेल्या तिमाहीत ८४ टक्के तक्रारी सोडवल्या गेल्या, तेव्हा ते खरे उत्तरदायित्व आहे. तेच नागरिक आणि सरकार यांच्यातील विश्वास निर्माण करते.

%% ── आमची टीम ────────────────────────────────────────────────────────────────
\section{आमची टीम}

\begin{tabularx}{\linewidth}{lX}
\toprule
\textbf{\textenglish{Manav Mutneja}} & सह-संस्थापक आणि मुख्य कार्यकारी अधिकारी. \textenglish{Antitrust} वकील आणि डिजिटल \textenglish{governance} तज्ज्ञ. धोरण, भागीदारी आणि सरकारी संस्थांशी संवाद यांचे नेतृत्व करतात. \\[4pt]
\textbf{\textenglish{Piyush Zaware}} & सह-संस्थापक आणि मुख्य तंत्रज्ञान अधिकारी. संगणकीय अभियंता आणि अर्थमितीशास्त्रज्ञ. तंत्रज्ञान पायाभूत सुविधा आणि \textenglish{data} प्रणाली तयार करतात. \\[4pt]
\textbf{\textenglish{Joe DiTommaso}} & सह-संस्थापक आणि मुख्य परिचालन अधिकारी. माहिती व्यवस्थापन आणि तंत्रज्ञान पार्श्वभूमी. उत्पादन परिचालन, वापरकर्ता संपादन आणि बाजारपेठेत प्रवेश यांचे नेतृत्व करतात. \\
\bottomrule
\end{tabularx}

\vspace{0.5em}
आम्ही सध्या \textenglish{Spring} २०२६ मध्ये थेट प्रात्यक्षिक \textenglish{demo}साठी तयारी करत आहोत.

%% ── आम्ही कुठे आहोत ──────────────────────────────────────────────────────────
\section{आम्ही कुठे आहोत?}

\begin{tcolorbox}[colback=lightgray, colframe=pollityblue, sharp corners, boxrule=0.5pt]
\begin{itemize}[leftmargin=1.5em, itemsep=3pt, topsep=2pt]
  \item[$\checkmark$] ग्राहक शोध मुलाखतींद्वारे मुख्य उत्पादन संकल्पना सिद्ध
  \item[$\checkmark$] जनश्रवण \textenglish{backend} तयार आणि तैनात
  \item[$\checkmark$] \textenglish{AI} वर्गीकरण प्रणाली कार्यरत (८ श्रेणी, ४ तातडी स्तर)
  \item[$\checkmark$] तक्रार \textenglish{Dashboard} तैनात
  \item[$\rightarrow$] \textenglish{WhatsApp Business} एकत्रीकरण प्रक्रियेत
  \item[$\rightarrow$] पहिले \textenglish{beta} वापरकर्ते लक्ष्य: ५० कार्यालये, तिसरी तिमाही २०२६
\end{itemize}
\end{tcolorbox}

\vspace{1em}
\begin{center}
{\color{pollityblue}\rule{0.5\linewidth}{0.4pt}}\\[0.5em]
{\small स्वारस्य आहे? संपर्क करा:}\\[0.3em]
{\small \textbf{\textenglish{piyush.zaware@pollity.in}} \textbar{} \textbf{\textenglish{manavmutneja@pollity.in}} \textbar{} \textbf{\textenglish{joe.ditommaso@pollity.in}}}
\end{center}

\end{document}
